% Article-ready fragment. Classical references and the source-report citation
% should be wired into the parent bibliography when integrating this file.
\section{Several quadratic factors beyond the first spectral derivative}
\label{msj:sec:main}

The incoming quadratic-spectrum report \cite{Quadratic} proves additivity of the first
spectral derivative for any even polynomial weight and any finite list
of masses. Its explicit next question asks for the second derivative,
the logarithmic cross term, and the limit when two masses coincide.
The following results answer that question with an ordinary convergent
series and a one-dimensional integral of elementary functions. They also
identify the finite corrections needed at every higher spectral order.
The classical multiplicative-anomaly framework \cite{EVZ} concerns first derivatives
of spectral zeta functions; the statements below concern the additional
higher derivatives. No historical-priority claim is needed for this
continuation of the inspected repository.

Fix $0<a<1$, put $v_n=(n+a)^2$ for $n\in\mathbb Z$, and let
\[
 P(v)=\sum_{d=0}^{D}p_dv^d,\qquad B_j>0\quad(1\le j\le r),\qquad r\ge2.
\]
The parameter $B_j$ is the squared mass. Initially in their respective
half-planes of absolute convergence, define
\begin{align}
 \mathcal E_P(s;a,B)&=\sum_{n\in\mathbb Z}P(v_n)(v_n+B)^{-s},
 \label{msj:eq:E}\\
 \mathcal Z_{P,\boldsymbol B}(s;a)
 &=\sum_{n\in\mathbb Z}P(v_n)
                  \prod_{j=1}^{r}(v_n+B_j)^{-s}.
 \label{msj:eq:Z}
\end{align}
Their continuations are holomorphic at $s=0$. Indeed, a tail expansion
uses integer powers of $v_n^{-1}$, and a possible pole at zero would
require an odd power $(n+a)^{-1}$. The detailed subtraction proof below
also supplies this continuation directly.

Write
\begin{equation}
 \mathfrak D_k(P;a,\boldsymbol B)
 =r^{k-1}\sum_{j=1}^{r}\mathcal E_P^{[k]}(0;a,B_j)
       -\mathcal Z_{P,\boldsymbol B}^{[k]}(0;a),
 \qquad k\ge1,
 \label{msj:eq:defect}
\end{equation}
where square brackets indicate derivatives in the spectral parameter.
The factors $r^{k-1}$ are part of the definition. They result from
comparing the factors at the same spectral scale:
\begin{equation}
 \mathfrak D(s)=\frac1r\sum_j\mathcal E_P(rs;a,B_j)
                           -\mathcal Z_{P,\boldsymbol B}(s;a),
 \qquad \mathfrak D_k=\mathfrak D^{[k]}(0).
 \label{msj:eq:common-scale}
\end{equation}
The already established first-derivative identity is $\mathfrak D_1=0$.

\subsection{An exact second-derivative formula}

For a function with a Laurent expansion at $v=\infty$, let
$\operatorname{Pol}_{\infty}$ retain its nonnegative integer powers of
$v$. Define
\begin{align}
 T_P(v;B,C)
 &=\operatorname{Pol}_{\infty}
       \left\{P(v)\log^2\frac{v+B}{v+C}\right\},
 \label{msj:eq:T2}\\
 S_P(a;B,C)
 &=\sum_{n\in\mathbb Z}
 \left\{P(v_n)\log^2\frac{v_n+B}{v_n+C}-T_P(v_n;B,C)\right\}.
 \label{msj:eq:S2}
\end{align}
Every term in the subtraction is explicit. In fact,
\begin{align}
 q_k(B,C)&=(-1)^k\sum_{m=1}^{k-1}
 \frac{(B^m-C^m)(B^{k-m}-C^{k-m})}{m(k-m)},\qquad k\ge2,
 \label{msj:eq:qcoeff}\\
 T_P(v;B,C)&=\sum_{d=2}^{D}p_d\sum_{k=2}^{d}q_k(B,C)v^{d-k}.
 \label{msj:eq:T2-explicit}
\end{align}
Empty sums are zero. Thus the weights $P(v)=1$ and $P(v)=v$ need no
subtraction. For $P(v)=v^2$ the only subtraction is $(B-C)^2$, while
for $P(v)=v^3$ it is
\[
 (B-C)^2v-(B-C)^2(B+C).
\]

\begin{theorem}[Second spectral derivative for any number of masses]
\label{msj:thm:second}
The series \eqref{msj:eq:S2} is ordinary and absolutely convergent, and
\begin{equation}
 \boxed{\displaystyle
 \mathfrak D_2(P;a,\boldsymbol B)
     =\sum_{1\le i<j\le r}S_P(a;B_i,B_j).}
 \label{msj:eq:second-pairwise}
\end{equation}
Consequently
\begin{equation}
 \mathcal Z_{P,\boldsymbol B}^{[2]}(0;a)
 =r\sum_j\mathcal E_P^{[2]}(0;a,B_j)
                       -\sum_{i<j}S_P(a;B_i,B_j).
 \label{msj:eq:second-evaluation}
\end{equation}
For $\deg P\le1$, the pair correction is the unsubtracted convergent
sum of logarithm squares. In particular, for $P=1$ it is nonzero whenever
at least two masses are distinct. Thus first-derivative additivity does
not extend unchanged to second derivatives.
\end{theorem}

\begin{proof}
At a fixed $v>0$ write $L_j=\log(v+B_j)$. Twice differentiating the
common-scale summand at zero gives
\[
 P(v)\left\{r\sum_jL_j^2-\left(\sum_jL_j\right)^2\right\}
        =P(v)\sum_{i<j}(L_i-L_j)^2.
\]
The logarithm difference has the expansion
\[
 \log\frac{v+B}{v+C}
   =\sum_{m\ge1}\frac{(-1)^{m+1}(B^m-C^m)}{m v^m}
\]
on a sufficiently large tail. Squaring proves
\eqref{msj:eq:qcoeff} and \eqref{msj:eq:T2-explicit}; the subtracted
summand is $O(v^{-1})$, uniformly for masses in a compact positive set.
It therefore has an absolutely convergent bilateral sum.

It remains to retain the exact value of the removed polynomial. Put
\begin{equation}
 \mathcal H_a(s)=\zeta(2s,a)+\zeta(2s,1-a).
 \label{msj:eq:H}
\end{equation}
For every integer $m\ge0$,
\begin{equation}
 \mathcal H_a(-m)=0,
 \label{msj:eq:trivial-zero}
\end{equation}
because $\zeta(-2m,x)=-B_{2m+1}(x)/(2m+1)$ and
$B_{2m+1}(1-a)=-B_{2m+1}(a)$. Hence each removed $v^m$ has zero
bilateral spectral value. The common-scale normal-convergence argument
in Theorem~\ref{msj:thm:alljets} makes this differentiation and restoration
precise and gives zero finite correction at order two. This proves
\eqref{msj:eq:second-pairwise}. With an affine weight, the raw tail is
$O(v^{-1})$ or better. With $P=1$ every pair summand is nonnegative and
a distinct pair has strictly positive summands, giving the final claim.
\end{proof}

For two masses, the mixed logarithmic coordinate is equivalently
\begin{equation}
 \operatorname{Reg}\sum_{n\in\mathbb Z}
  P(v_n)\log(v_n+B)\log(v_n+C)
 =\frac{\mathcal E_P^{[2]}(0;a,B)+\mathcal E_P^{[2]}(0;a,C)
                     -S_P(a;B,C)}2.
 \label{msj:eq:cross-term}
\end{equation}
Here $\operatorname{Reg}$ denotes the analytic continuation determined
by the common quadratic spectral coordinate. Formula
\eqref{msj:eq:S2}, or the ordinary integral below, makes the new
cross-term correction explicit; it does not subtract two unspecified
divergent sums.

\subsection{An elementary integral for the entire correction}

The classical bilateral resolvent, from cotangent partial fractions \cite{DLMF}, is
\begin{equation}
 R_a(t)=\sum_{n\in\mathbb Z}\frac1{v_n+t}
       =\frac{\pi}{\sqrt t}
        \frac{\sinh(2\pi\sqrt t)}
             {\cosh(2\pi\sqrt t)-\cos(2\pi a)},\qquad t>0.
 \label{msj:eq:resolvent}
\end{equation}
It follows by pairing the cotangent partial fractions at
$a+i\sqrt t$ and $a-i\sqrt t$. Set
\[
 f_P(t)=P(-t)R_a(t).
\]

\begin{theorem}[Ordinary logarithmic quadrature]
\label{msj:thm:quadrature}
For $B>C>0$,
\begin{equation}
 \boxed{\displaystyle
 S_P(a;B,C)=2\int_C^B f_P(t)
                     \log\frac{B-t}{t-C}\,\mathrm dt.}
 \label{msj:eq:quadrature}
\end{equation}
The integral is ordinary and absolutely convergent; its only endpoint
singularities are logarithmic. For arbitrary positive $B,C$, use
$S_P(a;B,C)=S_P(a;C,B)$ and $S_P(a;B,B)=0$. A form with a continuous
integrand on the square is
\begin{equation}
 S_P(a;B,C)=\int_C^B\int_C^B
      \frac{f_P(t)-f_P(u)}{u-t}\,\mathrm du\,\mathrm dt,
 \label{msj:eq:double-quadrature}
\end{equation}
where the diagonal value is $-f_P'(t)$.
\end{theorem}

\begin{proof}
Polynomial division gives
\[
 P(v)/(v+t)=Q_t(v)+P(-t)/(v+t)
\]
with a polynomial $Q_t$. Therefore subtraction of the polynomial part
and summation of the rational remainder give
\begin{align*}
 &\sum_{n\in\mathbb Z}(1-\operatorname{Pol}_{\infty})
          \frac{P(v_n)}{(v_n+t)(v_n+u)}\\
 &\hspace{35mm}=\frac{P(-t)R_a(t)-P(-u)R_a(u)}{u-t}.
\end{align*}
The statement at $u=t$ follows by continuity. For $(t,u)$ in a compact
positive square, every subtracted summand is uniformly $O(v_n^{-1})$.
Normal convergence permits integration of this identity over $[C,B]^2$.
Since
\[
 \int_C^B\frac{\mathrm dt}{v+t}=\log\frac{v+B}{v+C},
\]
and finite polynomial extraction commutes with these integrals, the
left side becomes \eqref{msj:eq:S2}. This proves
\eqref{msj:eq:double-quadrature}.

To reduce the square integral, omit temporarily the strip
$|u-t|\le\varepsilon$, split the numerator, and interchange $u,t$ in
its second term. The result tends to
\[
 2\int_C^B f_P(t)\bigl(\log(B-t)-\log(t-C)\bigr)\,\mathrm dt.
\]
The omitted-strip contribution tends to zero because the divided
difference extends continuously to the diagonal. The endpoint
logarithms are integrable. Thus the temporary symmetric exclusion is
only a proof device; the final integral needs no principal value or
finite-part convention.
\end{proof}

For example, at $a=1/2$, $P=1$, $B=4$, $C=1$, the exact correction is
\begin{equation}
 \sum_{n\in\mathbb Z}
 \log^2\frac{(n+1/2)^2+4}{(n+1/2)^2+1}
 =2\pi\int_1^4\frac{\tanh(\pi\sqrt t)}{\sqrt t}
                      \log\frac{4-t}{t-1}\,\mathrm dt.
 \label{msj:eq:example-unweighted}
\end{equation}
At the quartic weight $P(v)=v^2$, the corresponding exact identity is
\begin{equation}
 \sum_{n\in\mathbb Z}\left[
 (n+1/2)^4\log^2\frac{(n+1/2)^2+4}{(n+1/2)^2+1}-9\right]
 =2\pi\int_1^4 t^{3/2}\tanh(\pi\sqrt t)
                       \log\frac{4-t}{t-1}\,\mathrm dt.
 \label{msj:eq:example-weighted}
\end{equation}
Both series and both integrals are ordinary convergent expressions.

\subsection{Every confluent coefficient in finite resolvents}

Let $H_m=\sum_{j=1}^m1/j$ and $H_0=0$, and put
$\rho=\min(a,1-a)$. For $k\ge1$, define the rationally subtracted
resolvent moment
\begin{equation}
 R_{P,k}(a;C)
 =\sum_{n\in\mathbb Z}(1-\operatorname{Pol}_{\infty})
                         \frac{P(v_n)}{(v_n+C)^k}.
 \label{msj:eq:Rpk}
\end{equation}
As above, this is an absolutely convergent series after its explicit
polynomial division.

\begin{theorem}[Complete confluent expansion]
\label{msj:thm:confluent}
For $C>0$ and $|h|<C+\rho^2$,
\begin{align}
 S_P(a;C+h,C)
 &=2\sum_{k=2}^{\infty}\frac{(-1)^kH_{k-1}}{k}
                       h^k R_{P,k}(a;C),
 \label{msj:eq:confluent-resolvent}\\
 &=-2\sum_{k=2}^{\infty}\frac{H_{k-1}}{k!}
                       f_P^{(k-1)}(C)h^k.
 \label{msj:eq:confluent-elementary}
\end{align}
Here $S_P(a;C+h,C)$ denotes its principal-logarithm analytic
continuation from positive masses to this disk.  Indeed every
$v_n+C+h$ has positive real part on the disk, so the logarithms
and the explicitly subtracted series define a single holomorphic
function there.
Each coefficient has the finite reduction
\begin{equation}
 R_{P,k}(a;C)
 =\frac{(-1)^{k-1}}{(k-1)!}f_P^{(k-1)}(C)
 =\sum_{\ell=0}^{\min(D,k-1)}
       \frac{P^{(\ell)}(-C)}{\ell!}
                          \mathcal E_1(k-\ell;a,C).
 \label{msj:eq:Rpk-finite}
\end{equation}
The values $\mathcal E_1(N;a,C)$ at positive integers are obtained by
differentiating the elementary function \eqref{msj:eq:resolvent}, or
by the manuscript's finite polygamma partial fractions.
\end{theorem}

\begin{proof}
The elementary Taylor identity
\[
 \log^2(1+x)=2\sum_{k=2}^{\infty}\frac{(-1)^kH_{k-1}}{k}x^k,
 \qquad |x|<1,
\]
follows by differentiating both sides and using
$\log(1+x)/(1+x)=\sum_{m\ge1}(-1)^{m+1}H_mx^m$.
Apply it with $x=h/(v_n+C)$, extract the polynomial parts before
summation, and use locally uniform convergence on the stated disk.
This proves \eqref{msj:eq:confluent-resolvent}.

Differentiating $f_P(C)$ under the subtracted resolvent sum proves the
first equality in \eqref{msj:eq:Rpk-finite}. For the second, expand
$P(v)$ at $v=-C$. Terms with $\ell\ge k$ are polynomials and disappear
under $1-\operatorname{Pol}_{\infty}$; the other terms give the displayed
finite sum. The derivative formula yields
\eqref{msj:eq:confluent-elementary}. The function $f_P$ is meromorphic
in $C$, with possible poles only at $C=-v_n$. The nearest possible pole
is at distance $C+\rho^2$, which also proves the claimed disk of
convergence. Zeros of $P$ can only remove poles and enlarge this disk.
\end{proof}

The first terms are
\[
 S_P(a;C+h,C)=h^2R_{P,2}(a;C)-h^3R_{P,3}(a;C)
                   +\frac{11}{12}h^4R_{P,4}(a;C)+O(h^5).
\]
In particular, for $P=1$, $a=1/2$, and $C=1$,
\begin{equation}
 S_1(1/2;1+h,1)
 =\frac12\left(\pi\tanh\pi-\pi^2\operatorname{sech}^2\pi\right)h^2
                                                    +O(h^3).
 \label{msj:eq:confluent-example}
\end{equation}
This explicitly evaluates the first nonzero coefficient when two
quadratic factors merge.

\subsection{A finite correction formula at every spectral order}

The vanishing \eqref{msj:eq:trivial-zero} removes polynomial
counterterms at the second derivative. Higher derivatives also produce
polynomials times powers of $\log v$, whose spectral values need not
vanish. The following formula retains every such contribution.

Define
\begin{align}
 A(s,z)&=\frac1r\sum_{j=1}^r(1+B_jz)^{-rs}
                         -\prod_{j=1}^r(1+B_jz)^{-s}
          =\sum_{q\ge2}a_q(s)z^q,
 \label{msj:eq:A}\\
 a_q(s)&=(-1)^q\left\{
  \frac1r\sum_{j=1}^r\frac{(rs)_qB_j^q}{q!}
  -\sum_{q_1+\cdots+q_r=q}\prod_{j=1}^r
                      \frac{(s)_{q_j}B_j^{q_j}}{q_j!}\right\}.
 \label{msj:eq:aq}
\end{align}
The coefficient polynomials $a_q(s)$ are divisible by $s^2$; the
coefficients $a_0,a_1$ vanish identically. Put
\begin{equation}
 T(s,v)=\sum_{d=2}^{D}p_d\sum_{q=2}^{d}a_q(s)v^{d-q}.
 \label{msj:eq:T-spectral}
\end{equation}

\begin{theorem}[All-order ordinary remainder and finite Hurwitz jets]
\label{msj:thm:alljets}
For $s$ in a neighborhood of zero,
\begin{align}
 \mathfrak D(s)
 &=\sum_{n\in\mathbb Z}v_n^{-rs}
             \bigl(P(v_n)A(s,v_n^{-1})-T(s,v_n)\bigr)
 \notag\\
 &\quad+\sum_{d=2}^{D}p_d\sum_{q=2}^{d}
                           a_q(s)\mathcal H_a(rs+q-d).
 \label{msj:eq:regulator}
\end{align}
The first series is normally convergent near zero. For $k\ge2$, define
\begin{align}
 F_k(v)&=(-1)^kP(v)\left\{
     r^{k-1}\sum_j\log^k(v+B_j)
            -\left(\sum_j\log(v+B_j)\right)^k\right\},
 \label{msj:eq:raw-k}\\
 T_k(v)&=\sum_{d=2}^{D}p_d\sum_{q=2}^{d}
       \sum_{j=2}^{k}\binom{k}{j}a_q^{[j]}(0)
                (-r\log v)^{k-j}v^{d-q},
 \label{msj:eq:Tk}\\
 C_k&=\sum_{d=2}^{D}p_d\sum_{q=2}^{d}
       \sum_{j=2}^{k-1}\binom{k}{j}a_q^{[j]}(0)
                 r^{k-j}\mathcal H_a^{[k-j]}(q-d).
 \label{msj:eq:Ck}
\end{align}
Then the exact $k$th spectral defect is
\begin{equation}
 \boxed{\displaystyle
 \mathfrak D_k=\sum_{n\in\mathbb Z}\bigl(F_k(v_n)-T_k(v_n)\bigr)+C_k.}
 \label{msj:eq:alljets}
\end{equation}
The remainder series is ordinary and absolutely convergent. Its
summand is $O(v_n^{-1}(1+|\log v_n|^{k-2}))$. All correction terms are
finite Hurwitz-zeta jets at nonpositive even integers. For affine $P$,
both $T_k$ and $C_k$ vanish at every $k\ge2$.
\end{theorem}

\begin{proof}
At $s=0$, the two terms in \eqref{msj:eq:A} both equal one, and their
first $s$ derivatives both equal $-\sum_j\log(1+B_jz)$. Thus
$a_q(s)=O(s^2)$. Expansion to first degree in $z$ shows that $a_1=0$.
The Pochhammer expansion proves \eqref{msj:eq:aq} exactly.

On a sufficiently large lattice tail, Taylor expansion in $z=v^{-1}$
and the definition of $T$ give
\[
 P(v)A(s,v^{-1})-T(s,v)=O(v^{-1})
\]
uniformly for $s$ in a compact neighborhood of zero. Multiplication by
$v^{-rs}$ leaves a normally convergent series, for example on
$|\operatorname{Re}s|<1/(4r)$. Finitely many smaller $v_n$ cause no
difficulty: the exact functions $A(s,v_n^{-1})$ are used there,
without applying a possibly invalid small-$z$ expansion.

Restoring the removed monomials by
$\sum_n v_n^{d-q-rs}=\mathcal H_a(rs+q-d)$ gives
\eqref{msj:eq:regulator}, first in a common half-plane after a finite
head is separated and then by meromorphic continuation. Alternatively,
the same head-tail construction is a direct definition of this
continuation. None of the restored factors has a pole at zero:
$q-d$ is an integer, whereas the only pole of $\mathcal H_a$ is at
$1/2$. This proves the asserted holomorphy and normal convergence.

Differentiate \eqref{msj:eq:regulator} $k$ times. The first raw
derivative is \eqref{msj:eq:raw-k}, and the derivative of
$v^{-rs}T(s,v)$ is \eqref{msj:eq:Tk}. In the restored finite sum,
the terms with $j=0,1$ vanish because $a_q$ is divisible by $s^2$;
the term $j=k$ vanishes by \eqref{msj:eq:trivial-zero}. The remaining
terms are exactly \eqref{msj:eq:Ck}. Derivatives of the normal tail
give the stated polynomial-logarithmic bound. For $D\le1$, all finite
sums are empty.
\end{proof}

The second-derivative theorem follows immediately: at $k=2$ the
finite sum $C_2$ is empty, the raw derivative is the pairwise square,
and $T_2$ is precisely its polynomial part. At the third derivative
the correction can already be nonzero. For two masses and $P(v)=v^2$,
\begin{align}
 F_3(v)&=-3v^2\log\bigl((v+B)(v+C)\bigr)
                         \log^2\frac{v+B}{v+C},
 \label{msj:eq:third-raw}\\
 T_3(v)&=-6(B-C)^2\log v,
 \label{msj:eq:third-subtraction}\\
 C_3&=-12(B-C)^2\log(2\sin\pi a).
 \label{msj:eq:third-contact}
\end{align}
Indeed $a_2(s)=\tfrac12(B-C)^2s^2$ in this case, while Lerch's
identity and Gamma reflection give
\[
 \mathcal H_a'(0)=2\bigl(\zeta'(0,a)+\zeta'(0,1-a)\bigr)
                 =-2\log(2\sin\pi a).
\]
Formula \eqref{msj:eq:third-contact} is therefore a concrete check on
the sign and the scale factor in the general correction law. Discarding
it would change the third spectral derivative.

\subsection{Polylogarithmic coordinates for the finite corrections}

Define the symmetric polylogarithm
\[
 \mathcal C_s(a)=\frac{\operatorname{Li}_s(e^{2\pi ia})
                          +\operatorname{Li}_s(e^{-2\pi ia})}{2}.
\]
For real $s$ this is the real part of the first polylogarithm; the
symmetric definition retains holomorphic dependence on complex $s$.
The symmetric Hurwitz functional equation gives, for $m\ge0$ and
$u$ near zero,
\begin{equation}
 \mathcal H_a(-m+u)
 =4(-1)^m\frac{\Gamma(2m+1-2u)}{(2\pi)^{2m+1-2u}}
                  \sin(\pi u)\mathcal C_{2m+1-2u}(a).
 \label{msj:eq:H-polylog}
\end{equation}
Consequently $\mathcal H_a^{[h]}(-m)$ is a finite combination of
polylogarithm order derivatives at $2m+1$, of order at most $h-1$,
and Gamma derivatives at $2m+1$. The zero of $\sin(\pi u)$ is
responsible for this reduction in order. For example,
\begin{equation}
 \mathcal H_a'(-m)=
 \frac{2(-1)^m(2m)!}{(2\pi)^{2m}}\mathcal C_{2m+1}(a).
 \label{msj:eq:Hprime-polylog}
\end{equation}
Thus the finite corrections in \eqref{msj:eq:Ck} use polylogarithm
order derivatives only through order $k-3$. This is a finite reduction
of the correction, not a claim that the ordinary remainder in
\eqref{msj:eq:alljets} always reduces to finitely many polylogarithms.

\subsection{A confluent law simultaneously at every spectral order}

\begin{theorem}[All-order mass-coalescence law]
\label{msj:thm:allconfluent}
Let $B>0$ and fixed real numbers $c_1,\ldots,c_r$ be given, and set
$B_j=B+\varepsilon c_j$. For every $k\ge2$,
\begin{align}
 \mathfrak D_k(P;a,\boldsymbol B)
 &=\binom{k}{2}r^{k-2}\varepsilon^2
       \left(r\sum_jc_j^2-\Bigl(\sum_jc_j\Bigr)^2\right)
                    \mathcal E_P^{[k-2]}(2;a,B)
       +O(\varepsilon^3).
 \label{msj:eq:allconfluence}
\end{align}
The right side is well defined even when its defining raw series at
$s=2$ diverges: $\mathcal E_P(s;a,B)$ has no pole at an integer.
For $\deg P\le1$ it is an ordinary spectral derivative of a convergent
resolvent series. If all $c_j$ are equal, the masses remain identical
and every defect vanishes exactly.
\end{theorem}

\begin{proof}
For a fixed $v$, expand the common-scale expression in
\eqref{msj:eq:common-scale} through second order in $\varepsilon$.
The constant and linear coefficients cancel. The quadratic term is
\[
 \frac{s^2\varepsilon^2}{2}
 \left(r\sum_jc_j^2-\Bigl(\sum_jc_j\Bigr)^2\right)(v+B)^{-rs-2}.
\]
The same tail subtraction as above justifies this identity for the
continued families, locally uniformly in $s$ near zero. Hence
\[
 \mathfrak D(s)=\frac{s^2\varepsilon^2}{2}
 \left(r\sum_jc_j^2-\Bigl(\sum_jc_j\Bigr)^2\right)
                 \mathcal E_P(rs+2;a,B)+O(\varepsilon^3).
\]
Taking its $k$th derivative gives \eqref{msj:eq:allconfluence}.
The single-mass tail expansion has possible poles only at half-integers,
which proves the assertion at $s=2$. When all masses are equal,
$\mathcal Z_{P,\boldsymbol B}(s;a)=\mathcal E_P(rs;a,B)$, so
\eqref{msj:eq:common-scale} is identically zero.
\end{proof}

\subsection{A global Hurwitz-jet evaluation of the individual terms}

For completeness, the individual second derivatives in
\eqref{msj:eq:second-evaluation} can be evaluated without leaving an
unevaluated multivariate spectral zeta function. Choose an integer
$N\ge0$ such that $B<(N+\rho)^2$, and define
\[
 \mathcal H_{a,N}(s)=\zeta(2s,N+a)+\zeta(2s,N+1-a).
\]
Separating the $2N$ lattice terms with positive coordinates
$n+a$ and $n+1-a$, $0\le n<N$, gives the geometrically convergent
formula
\begin{align}
 \mathcal E_P^{[2]}(0;a,B)
 &=\sum_{n=0}^{N-1}\sum_{x\in\{n+a,n+1-a\}}
                         P(x^2)\log^2(x^2+B)
 \notag\\
 &\quad+\sum_{d=0}^{D}p_d\Bigg\{
    \mathcal H_{a,N}^{[2]}(-d)
 \notag\\
 &\hspace{20mm}+2\sum_{q=1}^{\infty}\frac{(-1)^qB^q}{q}
       \left[H_{q-1}\mathcal H_{a,N}(q-d)
                         +\mathcal H_{a,N}'(q-d)\right]\Bigg\}.
 \label{msj:eq:individual-second}
\end{align}
Every spectral argument here is an integer in the $\mathcal H$ variable,
so no pole is substituted separately. This identity follows directly
from
\[
 (s)_q=(q-1)!s\bigl(1+H_{q-1}s+O(s^2)\bigr)
\]
and the binomial expansion of $(x^2+B)^{-s}$ on the chosen tail.
The strict tail-radius condition justifies spectral differentiation
term by term. Higher derivatives use the same finite Pochhammer
coefficient calculus, or unsigned Stirling numbers of the first kind.
Together with \eqref{msj:eq:quadrature}, this supplies a global exact
evaluation of the several-mass second derivative in ordinary Hurwitz
jets and an explicitly anchored elementary integral.

\subsection{Verification and the remaining boundary}

The accompanying Python program constructs the Pochhammer spectral
coefficients and the direct logarithmic coefficients independently over
the rational numbers. It compares the resulting ordinary remainders
with independent finite-head Hurwitz continuations, evaluates the
resolvent quadrature directly, checks the nonzero third-order correction,
and compares the convergent confluent series with its integral.
Numerical diagnostics are distinct from the all-order proofs above;
they are not interval certificates.

The answered repository question is the several-mass second derivative
and its confluent limit, with the ordinary logarithmic cross term made
explicit. The one-dimensional integral in \eqref{msj:eq:quadrature}
is itself an exact evaluation coordinate. A finite reduction of that
integral to a specified algebra of ordinary polylogarithms for arbitrary
$a,B,C$ remains a separate question. For odd or noninteger spectral
weights, resonant odd powers can occur and additional local Laurent
terms must be retained; the parity-protected bilateral formula above
does not silently assert their absence.
